% RESULTADO_RECUPERADO. II REV10, 05_accion.tex:1--195,
% revision_planck.tex:13--93 y definicion_planck_rev06.tex:4--34.
% Selección: antecedentes de acción de la elipse; no incorpora aplicaciones
% electrónicas ni el contraste metrológico, que no determinan el radio.
\section{Structure déterminantale et sections orientées d’action}
\label{iv:action:section}

Les représentations de \(\pi,\varphi,e,\alpha\) proviennent de l’état
APP--TRIT--TPK des sections
\ref{sec:nucleo}--\ref{sec:alpha}. Le registre dodécaphasique de
\ref{sec:k-registro} conserve son orientation, ses résidus et
quotients et la mémoire du retour
\eqref{eq:holonomia-memoria}. Sur ce registre, on définit deux opérations
successives : une norme locale détermine la valuation de la section
d’action, et un caractère à continuation régionale détermine ses
coordonnées antérieure et postérieure au retour. Les deux opérations
transmettent à la construction angulaire leurs règles et leurs normalisations.

\subsection{Quotient local et norme déterminantale}
\label{iv:action:determinantal}

Le bloc local de pas trois du registre dodécaphasique, identifié
dans \ref{subsec:k-hadamard}, est
\begin{equation}
 H_4=\begin{pmatrix}
 1&1&1&1\\
 1&1&-1&-1\\
 1&-1&1&-1\\
 1&-1&-1&1
 \end{pmatrix},
 \qquad G_H=\mathbb Z^4/H_4\mathbb Z^4.
 \label{iv:action:h4}
\end{equation}
Le quotient correspond à ce bloc local. Les trois blocs du
registre complet conservent leurs positions ; l’opération de norme
qui suit agit sur les feuillets de \(G_H\).

\begin{lemma}[Structure du quotient local]
\label{iv:action:smith}
La forme normale de Smith de \(H_4\) est
\(\operatorname{diag}(1,2,2,4)\). Par conséquent,
\[
 G_H\simeq\mathbb Z/2\mathbb Z\oplus\mathbb Z/2\mathbb Z
       \oplus\mathbb Z/4\mathbb Z,\qquad |G_H|=16.
\]
\end{lemma}
\begin{proof}
Le plus grand commun diviseur des entrées est un. Les mineurs
d’ordre deux sont pairs et l’un a pour valeur \(-2\), de sorte que leur plus grand
commun diviseur est deux. Les identités
\(H_4H_4^{\mathsf T}=4I\) et \(\det H_4=-16\) montrent que les
cofacteurs valent \(\pm4\) ; le troisième diviseur déterminantal est
quatre. Le quatrième est \(16\). Les quotients successifs de
\(1,2,4,16\) sont \(1,2,2,4\) ; leur produit détermine l’ordre
du quotient.
\end{proof}

\begin{definition}[Lecteur déterminantal d'action]
\label{iv:action:norma}
La section scalaire \(\alpha\), constante sur les feuillets de \(G_H\),
et une application de l’auto-échelle \(\varphi\) déterminent
\begin{equation}
 \Sigma_H=\varphi\,\mathrm N_{G_H}(\alpha),\qquad
 \mathrm N_{G_H}(\alpha)=\prod_{g\in G_H}\alpha=\alpha^{16}.
 \label{iv:action:sigma}
\end{equation}
Son évaluation décimale est

\[
 \operatorname{ord}_{10}(\Sigma_H)
   =\lfloor\log_{10}\Sigma_H\rfloor.
\]
\end{definition}
La puissance seize provient du nombre de feuillets une fois fixée
la norme multiplicative. La norme et l’application unique de
\(\varphi\) font partie du lecteur déclaré : l’ordre du groupe,
considéré isolément, admet d’autres fonctionnelles.


\begin{theorem}[Valuation décimale de la section d’action]
\label{iv:action:decada}
Le lecteur précédent, appliqué aux coordonnées HMT déjà engendrées,
vérifie
\begin{equation}
 10^{-34}<\Sigma_H<10^{-33},
 \qquad \operatorname{ord}_{10}(\Sigma_H)=-34.
 \label{iv:action:orden-decimal}
\end{equation}
\end{theorem}
\begin{proof}
On utilise les inclusions rationnelles des représentations
précédentes,
\begin{equation}
 \frac{1618}{1000}<\varphi<\frac{1619}{1000},
 \qquad
 \frac{7297}{10^6}<\alpha<\frac{7298}{10^6}.
 \label{iv:action:intervalos}
\end{equation}
La fonction \((x,a)\mapsto xa^{16}\) est croissante en ses deux
variables positives. Les inégalités entières

\[
 1618\cdot7297^{16}>10^{65},
 \qquad 1619\cdot7298^{16}<10^{66}
\]
donnent
\[
 10^{-34}
 <\frac{1618}{1000}\left(\frac{7297}{10^6}\right)^{16}
 <\Sigma_H
 <\frac{1619}{1000}\left(\frac{7298}{10^6}\right)^{16}
 <10^{-33}.
\]
La définition de la valuation conclut la preuve. Les inclusions
sont des contrôles ultérieurs des coordonnées déjà produites ;
l’opération déterminantale précède toute mise en coordonnées des unités.
\end{proof}

\subsection{Caractère d’action et prolongement régional
}
\label{iv:action:caracter}

Le système de coordonnées central conserve la matrice
\[
 B_c=\begin{pmatrix}7&2\\2&7\end{pmatrix},
 \qquad\lambda_+=9,\quad\lambda_-=5.
\]
Le calendrier est de longueur \(12\), l’orbite du pas trois est de
longueur quatre et le quotient de recensement utilisé est
\(104976/1944=54\), avec les recensements définis dans la
section~\ref{sec:extension}. Ces invariants déterminent les coefficients
attribués aux degrés \(2,3,4,5\) du caractère :

\begin{equation}
 \frac9{16}=\frac{\lambda_+}{4^2},\qquad
 \frac59=\frac{\lambda_-}{\lambda_+},\qquad
 \frac7{48}=\frac{\operatorname{tr}(B_c)/2}{4\cdot12},
 \qquad
 \frac1{54}=\frac{1944}{104976}.
 \label{iv:action:coeficientes}
\end{equation}
L’attribution à ces degrés et l’alternance des signes appartiennent
à la définition du caractère analytique. Le contenu opératoire
réunit le support, l’orientation et cette règle d’attribution.

\begin{definition}[Caractère local et lecteur régional de retour]
\label{iv:action:lectores}
Sur les représentations \(\alpha,\varphi,\pi\), le caractère de
cinquième ordre, le transfert régional et sa continuation sont
\begin{align}
 H_5={}&\frac12\sqrt{\frac{1000\alpha}{\varphi}}-\alpha
       +\frac9{16}\alpha^2-\frac59\alpha^3
       +\frac7{48}\alpha^4-\frac1{54}\alpha^5,
       \label{iv:action:h5}\\
 D_A={}&\exp\!\left(-\frac{100\pi\alpha}{9}\right),
       \label{iv:action:da}\\
 \mathcal C_\pi(\alpha)={}&169\alpha^6+
       \frac{D_A\alpha^7}{1-D_A\alpha},
       \label{iv:action:cpi}\\
 \eta_{\mathrm{ret}}={}&H_5-\frac{90}{\pi}\mathcal C_\pi(\alpha).
       \label{iv:action:eta-ret}
\end{align}
L’entier \(169\) est le sélecteur régional des représentations
construites dans \ref{sec:generacion} : on conserve l’émission régionale
\(169\,|\,272\,|\,272\), distinguée de
\(418\,|\,674\,|\,521\) dans la construction exceptionnelle.
Le premier prolongement strict de sa continuation
a un poids \(D_A\) et chaque concaténation multiplie le poids par
\(D_A\). La normalisation initiale est unitaire.
\end{definition}

Le prolongement se formule avant d’évaluer en \(z=\alpha\).
Avec \(D_A\) déjà produit par son lecteur, nous écrivons
\(\mathscr C(z)=169z^6+\mathscr R(z)\), où
\(\mathscr R(z)\in z^7\mathbb R[[z]]\).
La concaténation augmente le degré d’une unité ; le poids initial
et la règle de transfert donnent

\begin{equation}
 \mathscr R(z)=D_Az^7+D_Az\mathscr R(z).
 \label{iv:action:renovacion}
\end{equation}

\begin{proposition}[Unicité de la continuation normalisée]
\label{iv:action:unicidad-renovacion}
L’équation \eqref{iv:action:renovacion} détermine une unique série :
\begin{equation}
 \mathscr C(z)=169z^6+\frac{D_Az^7}{1-D_Az}
             =169z^6+\sum_{n=7}^{\infty}D_A^{n-6}z^n.
 \label{iv:action:serie-renovacion}
\end{equation}
Sa réalisation est analytique pour \(|D_Az|<1\).
\end{proposition}
\begin{proof}
Le terme constant de \(1-D_Az\) est un, de sorte que cet
élément est inversible dans l’anneau des séries formelles.
Résoudre \eqref{iv:action:renovacion} donne l’expression rationnelle.
De manière équivalente, le coefficient de degré sept est \(D_A\) et chaque
coefficient ultérieur est \(D_A\) fois le précédent. La série
géométrique converge absolument dans le disque indiqué.
\end{proof}

La normalisation du premier poids intervient dans l’unicité.
L’impulsion de degré six et le rapport de concaténation, sans ce
poids initial, sont également compatibles avec

\[
 169z^6+\lambda\frac{D_Az^7}{1-D_Az}.
\]
La règle unitaire fixe \(\lambda=1\) avant l’évaluation.
Pour \(0<\alpha<1\), on a \(0<D_A<1\) et \(D_A\alpha<1\) ;
par conséquent,
\begin{equation}
 \frac{D_A\alpha^7}{1-D_A\alpha}
   =\sum_{n=0}^{\infty}D_A^{n+1}\alpha^{n+7}.
 \label{iv:action:cola}
\end{equation}
Après le terme d’indice \(N\), le reste exact est
\(D_A^{N+2}\alpha^{N+8}/(1-D_A\alpha)\). L’expression rationnelle
conserve ainsi la continuation complète à toute précision.

\subsection{Sections positives et action par tour
}
\label{iv:action:secciones}

Soit \(\mathcal L_S\) la droite orientée d’action et
\(\mathcal U_S\) une base positive. Ses sections antérieure et
postérieure au retour sont

\begin{equation}
 \hbar_{\mathrm{pre}}=H_5\,10^{-34}\mathcal U_S,\qquad
 \hbar_{\mathrm{ret}}=\eta_{\mathrm{ret}}\,10^{-34}\mathcal U_S.
 \label{iv:action:hbar}
\end{equation}
Le retour conserve la compensation régionale
\(90\mathcal C_\pi(\alpha)/\pi\) et le rapport
\begin{equation}
 \delta_{\mathrm{ret}}=H_5-\eta_{\mathrm{ret}},\qquad
 R_{\mathrm{act}}=\frac{\hbar_{\mathrm{pre}}}
                        {\hbar_{\mathrm{ret}}}
                =\frac{H_5}{\eta_{\mathrm{ret}}}.
 \label{iv:action:razon}
\end{equation}

\begin{proposition}[Positivité et covariance du rapport d’action]
\label{iv:action:positividad}
Dans les intervalles \eqref{iv:action:intervalos},
\(0<\eta_{\mathrm{ret}}<H_5\), et \(R_{\mathrm{act}}>1\).
Le rapport reste fixe lorsque la base positive commune
\(\mathcal U_S\) change.
\end{proposition}
\begin{proof}
Les termes de \(\mathcal C_\pi\) sont positifs. Pour borner
\(H_5\), \(\alpha>0.007\), \(\alpha<0.008\) et
\(\varphi<1.619\) suffisent. Le terme de racine de \eqref{iv:action:h5}
est supérieur à un et la somme de ses termes négatifs est inférieure
à \(0.008001\) ; donc \(H_5>0.99\).
Avec \(D_A<1\) et \(\pi>3\),
\[
 0<\frac{90}{\pi}\mathcal C_\pi
 <30\left(169(0.008)^6+
             \frac{(0.008)^7}{1-0.008}\right)<10^{-6}.
\]
Ainsi, \(\eta_{\mathrm{ret}}>0.989999>0\) et
\(\eta_{\mathrm{ret}}<H_5\). Le changement de base divise les deux
coordonnées d’action par le même facteur positif ; celui-ci
s’annule dans leur quotient.

\end{proof}

La détermination structurelle de l’action angulaire est la paire
\((\hbar_{\mathrm{pre}},\hbar_{\mathrm{ret}})\) avec son
identification par retour. Elle conserve la valuation de la norme
locale, le caractère d’action et la continuation régionale.
Dans chaque section \(s\in\{\mathrm{pre},\mathrm{ret}\}\), l’action
d’un tour complet est
\begin{equation}
 h_s=2\pi\hbar_s.
 \label{iv:action:vuelta}
\end{equation}
La mise en coordonnées ultérieure \(\mathcal U_S\mapsto\mathrm{J\,s}\) exprime
ces sections en unités d’action. Le rayon normalisé qui
sera construit conserve son indépendance par rapport à ce choix
commun d’unités.
